\documentclass{beamer}

\usetheme{Madrid}
\usecolortheme{default}
\usepackage[utf8]{inputenc}
\usepackage{graphicx} % Required for images

\title{Communication Protocols}
\subtitle{ME461 - Group: Little Daisies}
\author[Little Daisies]{Abdullah~Naci~Bodur\and
        Buğra~Çınaroğlu\and
        Ogan~Altuğ~Okutan
}
\date[ME461 Fall 25-26]{Fall 2025-2026}


\begin{document}

% ------------------------------------
\frame{\titlepage}
% ------------------------------------

% ------------------------------------
\section{Overview}
% ------------------------------------

\begin{frame}{Table of Contents}
\tableofcontents
\end{frame}

% ------------------------------------
\section{Introduction}
% ------------------------------------

\begin{frame}{Serial vs. Parallel}
\begin{itemize}
    \item \textbf{Parallel Communication:}
    \begin{itemize}
        \item Multiple wires carry multiple bits simultaneously.
        \item Fast but bulky and expensive (cables, crosstalk).
    \end{itemize}
    \item \textbf{Serial Communication:}
    \begin{itemize}
        \item Data is sent one bit at a time over a single channel.
        \item Reduces wiring complexity and cost.
        \item Standard for most modern embedded systems.
    \end{itemize}
\end{itemize}
\end{frame}

% ------------------------------------
\section{Asynchronous Communication (UART & RS-232)}
% ------------------------------------

\begin{frame}{Asynchronous Basics}
\begin{itemize}
    \item \textbf{No Clock Line:} Devices must agree on a timing speed (\textit{Baud Rate}) beforehand.
    \item \textbf{Synchronization:} Done via "Start" and "Stop" bits.
    \item \textbf{Common Protocols:} UART (Logic Level), RS-232 (Voltage Level).
\end{itemize}
\end{frame}

\begin{frame}{UART: Working Principle}
\textbf{Universal Asynchronous Receiver-Transmitter (UART)}
\begin{enumerate}
    \item \textbf{Parallel-to-Serial:} CPU writes byte to Transmit Register $\rightarrow$ Shift Register pushes bits out one by one.
    \item \textbf{Serial-to-Parallel:} Receiver samples line $\rightarrow$ shifts bits in $\rightarrow$ CPU reads byte.
\end{enumerate}

\begin{itemize}
    \item \textbf{Wiring:} Cross-coupled (TX connects to RX).
    \item \textbf{Topology:} Strictly Point-to-Point (1 Master, 1 Slave).
\end{itemize}
\end{frame}

\begin{frame}[fragile]{UART Frame Structure}
Standard 8N1 Frame (8 data bits, No parity, 1 stop bit):

\begin{verbatim}
Idle (High)
   |
   v
 Start | D0  D1  D2  D3  D4  D5  D6  D7 | Stop
  (0)  |            DATA                | (1)
\end{verbatim}

\begin{itemize}
    \item \textbf{Start Bit:} Pulls line LOW (0) to wake up receiver.
    \item \textbf{Data:} Sent LSB (Least Significant Bit) first.
    \item \textbf{Parity (Optional):} Error checking bit.
    \item \textbf{Stop Bit:} Pulls line HIGH (1) to reset for next byte.
\end{itemize}
\end{frame}

% --- ME461 EXAMPLE: UART ---
\begin{frame}{ME461 Application: UART Scope Project}
\begin{columns}
    \begin{column}{0.6\textwidth}
        \textbf{Project Context:}
        \begin{itemize}
            \item Used for \textbf{Raspberry Pi Pico programming} (REPL) and the \textbf{Scope Project}.
            \item The Pico reads signals/PWM and transmits data to the PC for visualization.
        \end{itemize}
        
        \vspace{0.2cm}
        \textbf{Configuration:}
        \begin{itemize}
            \item \textbf{Baud Rate:} 115200
            \item \textbf{Connection:} UART-over-USB or physical UART pins.
        \end{itemize}
    \end{column}
    
    \begin{column}{0.35\textwidth}
        \begin{block}{Pico UART Pinout}
        \centering
        \begin{tabular}{ll}
        \toprule
        \textbf{Function} & \textbf{Pin} \\
        \midrule
        \textbf{TX} & GP0 \\
        \textbf{RX} & GP1 \\
        \textbf{GND} & Any GND \\
        \bottomrule
        \end{tabular}
        \end{block}
        
        \vspace{0.2cm}
        \footnotesize \textit{*Scope project sends PWM values to Python script via this link.}
    \end{column}
\end{columns}
\end{frame}

\begin{frame}[fragile]{RS-232: The Physical Layer}
UART is a logic protocol; RS-232 is the voltage standard often used with it.
\begin{itemize}
    \item \textbf{Purpose:} Long distance communication (up to 15m) where 5V signals would degrade.
    \item \textbf{Voltage Levels (Inverted Logic):}
    \begin{itemize}
        \item Logic 0: +3V to +15V
        \item Logic 1: -3V to -15V
    \end{itemize}
\end{itemize}



\begin{verbatim}
          +15V  ------+       +-------
    Logic 0           |       |
          -15V        +-------+
    Logic 1
\end{verbatim}
\textit{Note: You need a MAX232 converter to connect this to a microcontroller.}
\end{frame}

% ------------------------------------
\section{Synchronous Communication (SPI)}
% ------------------------------------

\begin{frame}{SPI: Serial Peripheral Interface}
\textbf{Overview:}
\begin{itemize}
    \item Synchronous: Uses a dedicated clock line (SCK).
    \item Architecture: Master-Slave.
    \item Duplex: \textbf{Full Duplex} (Send and Receive simultaneously).
\end{itemize}

\textbf{Wiring (4 Lines):}
\begin{enumerate}
    \item \textbf{SCK:} Serial Clock (generated by Master).
    \item \textbf{MOSI:} Master Out Slave In (Data to slave).
    \item \textbf{MISO:} Master In Slave Out (Data to master).
    \item \textbf{CS/SS:} Chip Select (Active Low, selects the specific slave).
\end{enumerate}
\end{frame}

\begin{frame}{SPI Topology}

\begin{itemize}
    \item \textbf{Standard Mode:} One CS line per Slave.
    \item \textbf{Daisy Chain:} Slaves connected in series (Output of one goes to Input of next).
\end{itemize}

\textbf{Advantages:}
\begin{itemize}
    \item Fastest serial protocol (10MHz - 100MHz).
    \item No complex addressing (hardware selection via CS).
\end{itemize}
\textbf{Disadvantages:}
\begin{itemize}
    \item Requires many wires (3 + N slaves).
    \item No hardware acknowledgment (blind transmission).
\end{itemize}
\end{frame}

\begin{frame}[fragile]{SPI Timing Modes}
SPI has 4 modes based on Clock Polarity (CPOL) and Phase (CPHA).
\begin{itemize}
    \item \textbf{CPOL:} Idle state of clock (High or Low).
    \item \textbf{CPHA:} Data sampled on Rising or Falling edge.
\end{itemize}

\begin{verbatim}
Mode 0 (Standard):
SCK:   _   _   _   _
      / \_/ \_/ \_/ \_  (Sample on Rising Edge)
MOSI:  Valid Data
\end{verbatim}
\end{frame}

% --- ME461 EXAMPLE: SPI ---
\begin{frame}{ME461 Application: SPI Tetris Game}
\begin{columns}
    \begin{column}{0.6\textwidth}
        \textbf{Project Context:}
        \begin{itemize}
            \item Used in the \textbf{Tetris Game} project.
            \item High-speed communication required to update the LED grid matrix smoothly.
        \end{itemize}
        
        \vspace{0.2cm}
        \textbf{Hardware:}
        \begin{itemize}
            \item \textbf{Master:} Raspberry Pi Pico
            \item \textbf{Slave:} MAX7219 Dot Matrix Driver
        \end{itemize}
    \end{column}
    
    \begin{column}{0.35\textwidth}
        \begin{block}{Pico SPI0 Pinout}
        \centering
        \begin{tabular}{ll}
        \toprule
        \textbf{Function} & \textbf{Pin} \\
        \midrule
        \textbf{MOSI} & GP19 \\
        \textbf{SCK} & GP18 \\
        \textbf{CS (SS)} & GP17 \\
        \textbf{VCC} & 5V (VBUS) \\
        \bottomrule
        \end{tabular}
        \end{block}
        
        \vspace{0.2cm}
        \footnotesize \textit{*MISO is not used here (Display only receives data).}
    \end{column}
\end{columns}
\end{frame}




% ------------------------------------
\section{Synchronous Communication (I2C)}
% ------------------------------------

\begin{frame}{I\textsuperscript{2}C: Inter-Integrated Circuit}
\textbf{Overview:}
\begin{itemize}
    \item Synchronous, Half-Duplex.
    \item \textbf{2 Wires Only:} SDA (Data) and SCL (Clock).
    \item \textbf{Addressed Based:} No Chip Select lines; devices have a 7-bit address.
\end{itemize}

\textbf{Key Features:}
\begin{itemize}
    \item Multiple Masters and Multiple Slaves allowed.
    \item \textbf{Open Drain:} Requires Pull-up Resistors (4.7k$\Omega$).
    \item \textbf{Clock Stretching:} Slave can hold SCL low to pause Master.
\end{itemize}
\end{frame}

\begin{frame}[fragile]{I\textsuperscript{2}C Frame Structure}


[Image of I2C data packet structure]

\begin{verbatim}
[Start][Address (7bit)] [R/W] [ACK] [Data (8bit)] [ACK][Stop]
\end{verbatim}

\begin{enumerate}
    \item \textbf{Start Condition:} SDA goes LOW while SCL is HIGH.
    \item \textbf{Address:} Master broadcasts target address.
    \item \textbf{ACK/NACK:} Receiver pulls SDA low to acknowledge receipt.
    \item \textbf{Stop Condition:} SDA goes HIGH while SCL is HIGH.
\end{enumerate}
\end{frame}

\begin{frame}{I\textsuperscript{2}C Pros and Cons}
\textbf{Advantages:}
\begin{itemize}
    \item Low pin count (just 2 pins for up to 127 devices).
    \item Hardware flow control (ACK/NACK) ensures data integrity.
    \item Standard for sensors (IMUs, Temp sensors, EEPROMs).
\end{itemize}

\textbf{Disadvantages:}
\begin{itemize}
    \item Slower than SPI (Standard: 100kHz, Fast: 400kHz).
    \item Complex software implementation compared to SPI.
    \item Pull-up resistors consume power.
\end{itemize}
\end{frame}

% --- ME461 EXAMPLE: I2C ---
\begin{frame}{ME461 Application: I2C Mechaboard}
\begin{columns}
    \begin{column}{0.6\textwidth}
        \textbf{Project Context:}
        \begin{itemize}
            \item Used on the \textbf{Mechaboard} platform.
            \item \textbf{Bus Sharing:} Two different sensors are controlled using the same single pair of I2C pins.
        \end{itemize}
        
        \vspace{0.2cm}
        \textbf{Devices on Bus:}
        \begin{enumerate}
            \item \textbf{MPU6050:} IMU (Accel/Gyro)
            \item \textbf{SSD1306:} OLED Display
        \end{enumerate}
    \end{column}
    
    \begin{column}{0.35\textwidth}
        \begin{block}{Pico I2C0 Pinout}
        \centering
        \begin{tabular}{ll}
        \toprule
        \textbf{Function} & \textbf{Pin} \\
        \midrule
        \textbf{SDA} & GP4 \\
        \textbf{SCL} & GP5 \\
        \textbf{VCC} & 3.3V \\
        \bottomrule
        \end{tabular}
        \end{block}
        
        \vspace{0.2cm}
        \footnotesize \textit{*Different addresses allow both chips to talk on GP4/GP5 simultaneously.}
    \end{column}
\end{columns}
\end{frame}




% ------------------------------------
\section{Summary & Comparison}
% ------------------------------------

\begin{frame}{Comparison Matrix}
\begin{center}
\resizebox{\textwidth}{!}{%
\begin{tabular}{|l|c|c|c|}
\hline
\textbf{Feature} & \textbf{UART / RS-232} & \textbf{SPI} & \textbf{I\textsuperscript{2}C} \\
\hline
\textbf{Wires} & 2 (TX, RX) & 4 (SCK, MOSI, MISO, CS) & 2 (SDA, SCL) \\
\hline
\textbf{Synchronous} & No (Async) & Yes & Yes \\
\hline
\textbf{Speed} & Slow (<115 kbps typical) & Fastest (>10 Mbps) & Medium (100-400 kbps) \\
\hline
\textbf{Topology} & Point-to-Point & Single Master, Multi-Slave & Multi-Master, Multi-Slave \\
\hline
\textbf{Flow Control} & None / Software & None & ACK/NACK \\
\hline
\textbf{Duplex} & Full Duplex & Full Duplex & Half Duplex \\
\hline
\textbf{Best For} & PC $\leftrightarrow$ MCU, Debugging & Screens, SD Cards & Sensors, Low Pin Count \\
\hline
\end{tabular}%
}
\end{center}
\end{frame}

\begin{frame}{Conclusion}
\textbf{How to choose?}
\begin{itemize}
    \item Need simple PC logging or long distance? $\rightarrow$ \textbf{UART / RS-232}
    \item Need high speed data transfer (Display, Flash Memory)? $\rightarrow$ \textbf{SPI}
    \item Need many sensors with few pins? $\rightarrow$ \textbf{I\textsuperscript{2}C}
\end{itemize}

\vspace{1cm}
\begin{center}
\textbf{Questions?}
\end{center}
\end{frame}

\end{document}